\documentclass{beamer}

% Theme choice
\usetheme{Madrid} % You can change to e.g., Warsaw, Berlin, CambridgeUS, etc.

% Encoding and font
\usepackage[utf8]{inputenc}
\usepackage[T1]{fontenc}

% Graphics and tables
\usepackage{graphicx}
\usepackage{booktabs}

% Code listings
\usepackage{listings}
\lstset{
basicstyle=\ttfamily\small,
keywordstyle=\color{blue},
commentstyle=\color{gray},
stringstyle=\color{red},
breaklines=true,
frame=single
}

% Math packages
\usepackage{amsmath}
\usepackage{amssymb}

% Colors
\usepackage{xcolor}

% TikZ and PGFPlots
\usepackage{tikz}
\usepackage{pgfplots}
\pgfplotsset{compat=1.18}
\usetikzlibrary{positioning}

% Hyperlinks
\usepackage{hyperref}

% Title information
\title{Chapter 12: Capstone Project Development}
\author{Your Name}
\institute{Your Institution}
\date{\today}

\begin{document}

\frame{\titlepage}

\begin{frame}[fragile]
    \frametitle{Introduction to Capstone Project Development}
    \begin{block}{Overview of the Capstone Project}
        A \textbf{Capstone Project} is a culminating academic experience that allows students to apply their knowledge and skills in a real-world context. In the Python Fundamentals course, the capstone serves as a key opportunity for students to integrate what they have learned throughout the course.
    \end{block}
\end{frame}

\begin{frame}[fragile]
    \frametitle{Significance in the Python Fundamentals Course}
    \begin{itemize}
        \item \textbf{Practical Application}: Solidifies understanding of Python programming.
        \item \textbf{Comprehensive Skill Demonstration}: Showcases proficiency in programming logic, data manipulation, and library utilization.
        \item \textbf{Real-World Problem Solving}: Develops critical thinking and problem-solving skills.
    \end{itemize}
\end{frame}

\begin{frame}[fragile]
    \frametitle{Objectives for Student Learning}
    \begin{enumerate}
        \item \textbf{Problem Identification}: Learn to identify relevant problems and create solutions using Python.
        \begin{block}{Example}
            A student may choose to create a program that automates the sorting of large datasets for better efficiency.
        \end{block}
        
        \item \textbf{Application of Python Skills}: Encapsulates data structures, algorithms, and file management.
        \begin{lstlisting}[language=Python]
# Example of reading and processing data from a file
with open('data.txt', 'r') as file:
    data = file.read()
    lines = data.split('\n')
            \end{lstlisting}
        
        \item \textbf{Collaboration}: Encourages teamwork and effective communication.
    \end{enumerate}
\end{frame}

\begin{frame}[fragile]
    \frametitle{Key Points to Emphasize}
    \begin{itemize}
        \item \textbf{Integration of Knowledge}: The capstone bridges classroom learning and practical application.
        \item \textbf{Project Management Skills}: Experience in planning, executing, and finalizing projects.
        \item \textbf{Feedback Loop}: Regular check-ins encourage iterative refinement of projects.
    \end{itemize}
\end{frame}

\begin{frame}[fragile]
    \frametitle{Conclusion}
    The capstone project is a vital demonstration of individual learning and a stepping stone to a career in technology. Understanding its significance and objectives empowers students to approach this project with creativity and confidence, leading to impactful results.
\end{frame}

\begin{frame}[fragile]
    \frametitle{Project Objectives - Overview}
    \begin{itemize}
        \item Capstone project as the culmination of Python Fundamentals course.
        \item Application of learning to solve real-world problems.
        \item Emphasis on Python programming and collaboration.
    \end{itemize}
\end{frame}

\begin{frame}[fragile]
    \frametitle{Project Objectives - Problem Identification}
    \begin{block}{1. Problem Identification}
        \begin{itemize}
            \item \textbf{Definition}: Identify a relevant and meaningful problem to solve.
            \item \textbf{Approach}:
                \begin{itemize}
                    \item Conduct research in domains like healthcare, finance, education.
                    \item Engage in brainstorming sessions for idea generation.
                \end{itemize}
            \item \textbf{Example}: Streamlining appointment scheduling for a healthcare clinic using Python solutions.
        \end{itemize}
    \end{block}
\end{frame}

\begin{frame}[fragile]
    \frametitle{Project Objectives - Application of Python Skills}
    \begin{block}{2. Application of Python Skills}
        \begin{itemize}
            \item \textbf{Definition}: Apply various Python skills to create functional solutions.
            \item \textbf{Key Python Concepts}:
                \begin{itemize}
                    \item Data structures (lists, dictionaries, sets)
                    \item Control flow (if statements, loops)
                    \item Functions and modules for organization
                    \item Libraries (e.g., pandas, Flask)
                \end{itemize}
            \item \textbf{Example}: Using pandas to analyze data and visualize trends related to the identified problem.
        \end{itemize}
    \end{block}
\end{frame}

\begin{frame}[fragile]
    \frametitle{Project Objectives - Collaboration}
    \begin{block}{3. Collaboration}
        \begin{itemize}
            \item \textbf{Definition}: Work effectively with peers for enhanced learning.
            \item \textbf{Approach}:
                \begin{itemize}
                    \item Form small teams for communication and teamwork.
                    \item Assign roles based on strengths (e.g., programmer, analyst, manager).
                \end{itemize}
            \item \textbf{Example}: Collaborating with a data visualization expert for impactful presentations.
        \end{itemize}
    \end{block}
\end{frame}

\begin{frame}[fragile]
    \frametitle{Key Points and Conclusion}
    \begin{itemize}
        \item Objectives are interconnected; effective problem identification supports Python skills and collaboration.
        \item Encourage innovation and creativity in problem-solving.
        \item Focus on clear communication in both coding and presentations.
    \end{itemize}
    
    \begin{block}{Conclusion}
        By achieving these objectives, students solidify Python knowledge and develop critical problem-solving and teamwork skills for future careers.
    \end{block}
\end{frame}

\begin{frame}[fragile]
    \frametitle{Project Proposal Stage - Introduction}
    \begin{block}{Overview}
        The Project Proposal Stage is a critical point in your capstone project development. 
        It sets the foundation for your work by clearly articulating your project idea, objectives, and the approach you will take.
    \end{block}
\end{frame}

\begin{frame}[fragile]
    \frametitle{Project Proposal Stage - Proposal Requirements}
    \begin{block}{1. Format}
        The proposal should be formatted as a formal document with the following sections:
        \begin{itemize}
            \item \textbf{Title Page:} Include project title, your name, course details, and submission date.
            \item \textbf{Abstract:} A brief summary (150-200 words) outlining the essence of your proposal.
            \item \textbf{Introduction:} State the problem or need your project addresses.
            \item \textbf{Objectives:} Clearly define what you aim to accomplish.
            \item \textbf{Methodology:} Describe tools and techniques (e.g., Python libraries).
            \item \textbf{Timeline:} Outline key phases and estimated deadlines.
            \item \textbf{Feasibility:} Discuss the viability of your project.
            \item \textbf{References:} Include any prior work or resources referenced.
        \end{itemize}
    \end{block}
\end{frame}

\begin{frame}[fragile]
    \frametitle{Project Proposal Stage - Proposal Example and Timeline}
    \begin{block}{Example of Format}
    \begin{verbatim}
    Title: Enhancing Data Visualization in Python
    Author: [Your Name]
    Course: Capstone Project - Data Science 101
    Date: [Submission Date]

    Abstract: This project aims to create an interactive dashboard that visualizes real-time data...
    \end{verbatim}
    \end{block}

    \begin{block}{2. Content}
        \begin{itemize}
            \item \textbf{Clarity:} Be precise and concise. Avoid jargon unless defined.
            \item \textbf{Feasibility:} Assess whether your project can be completed within the given time frame.
        \end{itemize}
    \end{block}
    
    \begin{block}{3. Timeline}
        Ensure a realistic timeline, consider using a Gantt chart, and submit by the assigned deadline.
    \end{block}
\end{frame}

\begin{frame}[fragile]
    \frametitle{Project Proposal Stage - Key Points and Tips}
    \begin{block}{Key Points to Emphasize}
        \begin{itemize}
            \item \textbf{Be Clear and Concise:} Clarity in writing ensures understanding.
            \item \textbf{Feasibility is Crucial:} Choose a project scope that matches your skill level.
            \item \textbf{Follow the Structure:} Enhances professionalism and helps navigation.
        \end{itemize}
    \end{block}

    \begin{block}{Tips for Success}
        \begin{itemize}
            \item Draft multiple versions and seek peer review.
            \item Use clear visuals like charts and graphs.
            \item Utilize available resources to improve your document.
        \end{itemize}
    \end{block}
\end{frame}

\begin{frame}[fragile]
    \frametitle{Project Proposal Stage - Conclusion}
    \begin{block}{Conclusion}
        A well-crafted project proposal serves as a roadmap for your capstone project. It effectively conveys your intentions and sets you up for successful development. 
        Take your time to refine your proposal based on feedback!
    \end{block}
\end{frame}

\begin{frame}[fragile]
    \frametitle{Progress Report and Milestones - Overview}
    A progress report is a crucial tool in managing your capstone project, documenting advancements, challenges, and next steps. It serves two main purposes: 
    \begin{itemize}
        \item Formal accountability exercise
        \item Reflection on work during the development phase
    \end{itemize}
    
    \begin{block}{Key Elements of a Progress Report}
        \begin{itemize}
            \item \textbf{Project Status}: Summarize overall progress relative to the project timeline.
            \item \textbf{Completed Tasks}: List major accomplishments since the last report.
            \item \textbf{Upcoming Tasks}: Outline immediate goals and deliverables for the next period.
            \item \textbf{Challenges and Solutions}: Identify obstacles encountered and solutions implemented.
            \item \textbf{Timeline Updates}: Adjust the project timeline based on any delays or advancements.
        \end{itemize}
    \end{block}
\end{frame}

\begin{frame}[fragile]
    \frametitle{Progress Report and Milestones - Milestones}
    Milestones indicate significant points in your project timeline, marking the completion of key phases or deliverables. They help gauge project health and progress. 

    \begin{block}{Key Milestones Examples}
        \begin{enumerate}
            \item \textbf{Initial Research Complete}: Finish gathering background material and identifying stakeholder needs.
            \item \textbf{Prototype Development}: Completion of the first iteration of your project, ready for testing.
            \item \textbf{User Testing}: Conduct evaluations with actual users to receive feedback.
            \item \textbf{Final Submission Ready}: Ensure all aspects of the project are finalized and documented for submission.
        \end{enumerate}
    \end{block}
\end{frame}

\begin{frame}[fragile]
    \frametitle{Progress Report and Milestones - Tracking Progress}
    To ensure you stay on track, employ strategies for consistent monitoring and evaluation:

    \begin{itemize}
        \item \textbf{Gantt Chart}: Create a visual timeline including tasks, milestones, and deadlines.
        \begin{lstlisting}
Task Name      | Start Date | End Date | Status
----------------------------------------------------
Research       | Jan 1     | Jan 15   | Completed
Prototype Dev  | Jan 16    | Feb 15   | In Progress
Testing        | Feb 16    | Mar 1    | Pending
        \end{lstlisting}
        
        \item \textbf{Weekly Check-ins}: Schedule regular meetings with your advisor or peers.
        \item \textbf{Reflective Journals}: Maintain a project journal to document achievements and setbacks.
    \end{itemize}

    \begin{block}{Conclusion}
        Progress reports and milestones are crucial for your capstone project’s success. 
        Regular tracking enhances accountability and promotes effective project management.
    \end{block}
\end{frame}

\begin{frame}[fragile]
    \frametitle{Development Phase - Overview}
    \begin{block}{Overview of the Development Phase}
        The Development Phase of your capstone project is crucial where your ideas take shape. 
        This includes actual creation of the prototype through coding, testing, and iterating based on feedback.
    \end{block}
\end{frame}

\begin{frame}[fragile]
    \frametitle{Development Phase - Key Concepts}
    \begin{enumerate}
        \item \textbf{Prototyping}:
        \begin{itemize}
            \item Build a functioning version of your final product.
            \item A prototype is a preliminary model for demonstrating ideas and gathering feedback.
            \item Encourage rapid prototyping for iterative improvement.
        \end{itemize}
        
        \item \textbf{Coding Standards}:
        \begin{itemize}
            \item Follow industry-standard practices for readability.
            \item Important practices:
            \begin{itemize}
                \item Comment your code.
                \item Use consistent naming conventions.
                \item Utilize version control (e.g., Git).
            \end{itemize}
        \end{itemize}

        \item \textbf{Testing}:
        \begin{itemize}
            \item Implement unit and integration tests.
            \item Maintain a feedback loop for continuous improvement.
        \end{itemize}
    \end{enumerate}
\end{frame}

\begin{frame}[fragile]
    \frametitle{Development Phase - Expectations and Coding Reviews}
    \begin{block}{Expectations for Final Prototypes}
        \begin{itemize}
            \item \textbf{Functionality}: Demonstrate core functionality.
            \item \textbf{User Experience (UX)}: Design for intuitive navigation.
            \item \textbf{Documentation}: Outline system architecture, installation, and user manuals.
        \end{itemize}
    \end{block}

    \begin{block}{Coding Reviews}
        \begin{itemize}
            \item Participate in peer coding reviews.
            \item Focus areas:
            \begin{itemize}
                \item Clarity of the code.
                \item Performance and efficiencies.
                \item Bug identification.
            \end{itemize}
            \item Example Review Questions:
            \begin{itemize}
                \item Does the code adhere to standards?
                \item How could this function be optimized?
            \end{itemize}
        \end{itemize}
    \end{block}
\end{frame}

\begin{frame}[fragile]
    \frametitle{Final Presentation}
    \begin{block}{Description}
        Guidelines for the Final Presentation
    \end{block}
\end{frame}

\begin{frame}[fragile]
    \frametitle{Overview}
    \begin{itemize}
        \item The final presentation showcases your capstone project, skills, and results.
        \item Serves as a summary of your project journey.
        \item Aims to provide clarity about your work and its significance.
    \end{itemize}
\end{frame}

\begin{frame}[fragile]
    \frametitle{Required Elements}
    \begin{enumerate}
        \item \textbf{Introduction}
            \begin{itemize}
                \item Introduce yourself and your project.
                \item State the problem addressed and its significance.
            \end{itemize}
        \item \textbf{Objectives}
            \begin{itemize}
                \item Outline the goals and objectives of your project.
                \item Define what you aim to achieve.
            \end{itemize}
        \item \textbf{Methodology}
            \begin{itemize}
                \item Describe the methods and approaches used.
                \item Highlight essential tools, frameworks, or technologies.
            \end{itemize}
        \item \textbf{Results}
            \begin{itemize}
                \item Present project outcomes with visuals (graphs, charts).
            \end{itemize}
        \item \textbf{Challenges Faced}
            \begin{itemize}
                \item Discuss obstacles encountered and solutions.
            \end{itemize}
        \item \textbf{Conclusion}
            \begin{itemize}
                \item Summarize findings and implications.
                \item Offer recommendations for future work.
            \end{itemize}
        \item \textbf{Q\&A Session}
            \begin{itemize}
                \item Prepare for audience questions.
                \item Anticipate common queries.
            \end{itemize}
    \end{enumerate}
\end{frame}

\begin{frame}[fragile]
    \frametitle{Presentation Structure}
    \begin{itemize}
        \item \textbf{Duration}: 10-15 minutes.
        \item \textbf{Format}: Use PowerPoint or similar tool.
        \item \textbf{Clarity}: Keep the narrative clear and concise; avoid jargon.
        \item \textbf{Visual Aids}: Support spoken words with slides containing key points.
    \end{itemize}
\end{frame}

\begin{frame}[fragile]
    \frametitle{Assessment Criteria}
    \begin{enumerate}
        \item \textbf{Content Quality (40\%)}
        \item \textbf{Presentation Style (30\%)}
        \item \textbf{Q\&A Effectiveness (20\%)}
        \item \textbf{Time Management (10\%)}
    \end{enumerate}
\end{frame}

\begin{frame}[fragile]
    \frametitle{Key Points to Emphasize}
    \begin{itemize}
        \item Practice multiple times for fluency and timing.
        \item Use note cards or bullet points for guidance.
        \item Engage with the audience via eye contact and interaction.
        \item Ensure all technology works prior to the presentation.
    \end{itemize}
\end{frame}

\begin{frame}[fragile]
    \frametitle{Example Slide Layout}
    \begin{enumerate}
        \item Slide 1: Title Slide
        \item Slide 2: Introduction \& Objectives
        \item Slide 3: Methodology Overview
        \item Slide 4: Results with Key Data
        \item Slide 5: Challenges and Solutions
        \item Slide 6: Conclusion and Q\&A
    \end{enumerate}
\end{frame}

\begin{frame}[fragile]
    \frametitle{Grading Rubric - Overview}
    \begin{block}{Understanding the Grading Rubric}
        The grading rubric for the capstone project assesses the quality and effectiveness of your work based on specific criteria.
    \end{block}
    \begin{itemize}
        \item Components of the rubric:
            \begin{itemize}
                \item Functionality
                \item Complexity
                \item Presentation Quality
                \item Collaboration and Teamwork
            \end{itemize}
    \end{itemize}
\end{frame}

\begin{frame}[fragile]
    \frametitle{Grading Rubric - Evaluation Criteria}
    \begin{enumerate}
        \item \textbf{Functionality (30 points)}
            \begin{itemize}
                \item Completeness of features
                \item User interaction and experience
                \item Error handling and stability
            \end{itemize}
        \item \textbf{Complexity (30 points)}
            \begin{itemize}
                \item Innovative use of technology
                \item Difficulty level
                \item Depth of features
            \end{itemize}
        \item \textbf{Presentation Quality (20 points)}
            \begin{itemize}
                \item Clarity of presentation
                \item Engagement with the audience
                \item Professionalism
            \end{itemize}
        \item \textbf{Collaboration and Teamwork (20 points)}
            \begin{itemize}
                \item Equal participation
                \item Effective communication within the team
                \item Peer evaluations
            \end{itemize}
    \end{enumerate}
\end{frame}

\begin{frame}[fragile]
    \frametitle{Grading Rubric - Key Points}
    \begin{itemize}
        \item Familiarize yourself with the rubric early in the project.
        \item Regularly check against criteria to ensure you are on track.
        \item Provide evidence for functionality and teamwork sections.
        \item Effective presentations reflect both project quality and clarity.
    \end{itemize}
    
    \begin{block}{Conclusion}
        Understanding the grading rubric is crucial for success. Use it as a roadmap to guide your development ensuring a well-rounded and high-quality presentation.
    \end{block}
\end{frame}

\begin{frame}[fragile]
    \frametitle{Collaboration and Teamwork - Overview}
    \begin{block}{Importance of Collaboration}
        Collaboration is a key aspect of the capstone project that:
        \begin{itemize}
            \item Enhances creativity
            \item Fosters critical thinking
            \item Utilizes unique skills from each member
        \end{itemize}
        In effective teams, collaboration generates synergy where the total output is greater than the sum of individual efforts.
    \end{block}
\end{frame}

\begin{frame}[fragile]
    \frametitle{Collaboration and Teamwork - Roles and Responsibilities}
    \begin{block}{Roles and Responsibilities}
        Clear roles are essential for team effectiveness. Here are some typical roles:
        \begin{itemize}
            \item \textbf{Project Manager}: Oversees timeline and facilitates meetings.
            \item \textbf{Research Lead}: Gathers and synthesizes information.
            \item \textbf{Technical Developer}: Builds and codes the project.
            \item \textbf{Designer}: Focuses on visual and user interface aspects.
        \end{itemize}
        Early discussions and documentation of these roles can prevent overlap and ensure accountability.
    \end{block}
\end{frame}

\begin{frame}[fragile]
    \frametitle{Collaboration and Teamwork - Peer Evaluation and Responsibilities}
    \begin{block}{Peer Evaluations and Collective Responsibilities}
        \begin{itemize}
            \item Regular evaluations provide constructive feedback.
            \item Implement a 360-degree feedback system for:
                \begin{itemize}
                    \item Mid-project reviews to assess progress
                    \item Final evaluations reflecting on contributions
                \end{itemize}
        \end{itemize}
        
        \begin{block}{Collective Responsibilities}
            A successful team:
            \begin{itemize}
                \item Engages in key discussions and decisions
                \item Holds each other accountable for deadlines
                \item Distributes workload evenly
            \end{itemize}
            Key points to emphasize:
            \begin{itemize}
                \item Open communication
                \item Conflict resolution strategies
                \item Assessment of teamwork through tools like surveys
            \end{itemize}
        \end{block}
    \end{block}
\end{frame}

\begin{frame}[fragile]
    \frametitle{Feedback Mechanisms - Overview}
    \begin{block}{Understanding Feedback in the Capstone Project}
        Feedback is essential throughout the capstone project development process. It provides guidance, enhances learning, and helps ensure the project meets the expected standards. Effective feedback comes from various sources at multiple stages.
    \end{block}
\end{frame}

\begin{frame}[fragile]
    \frametitle{Feedback Mechanisms - Types of Feedback}
    \begin{enumerate}
        \item Instructor Feedback
        \item Peer Reviews
        \item Self-Reflection Practices
    \end{enumerate}
\end{frame}

\begin{frame}[fragile]
    \frametitle{Instructor Feedback}
    \begin{itemize}
        \item \textbf{Purpose:} Guides improvement with formative and summative feedback.
        \item \textbf{Examples:}
            \begin{itemize}
                \item Weekly Check-ins: Regular meetings to discuss progress.
                \item Draft Submissions: Submitting drafts for structured critique.
            \end{itemize}
        \item \textbf{Benefits:}
            \begin{itemize}
                \item Access to expert advice and insights.
                \item Early identification of potential pitfalls.
            \end{itemize}
    \end{itemize}
\end{frame}

\begin{frame}[fragile]
    \frametitle{Peer Reviews}
    \begin{itemize}
        \item \textbf{Purpose:} Allows for diverse perspectives and constructive criticism.
        \item \textbf{Examples:}
            \begin{itemize}
                \item Structured Peer Review Sessions: Guidelines help critical reviews.
                \item Feedback Forums: Online platforms for posting work and receiving input.
            \end{itemize}
        \item \textbf{Benefits:}
            \begin{itemize}
                \item Promotes collaborative learning.
                \item Encourages accountability and communication skills.
            \end{itemize}
    \end{itemize}
\end{frame}

\begin{frame}[fragile]
    \frametitle{Self-Reflection Practices}
    \begin{itemize}
        \item \textbf{Purpose:} Identifies strengths and areas for improvement.
        \item \textbf{Examples:}
            \begin{itemize}
                \item Reflection Journals: Document thoughts, challenges, and learnings.
                \item Guided Reflection Questions: Addressing "What went well?", "What didn't?", and "How can I improve?".
            \end{itemize}
        \item \textbf{Benefits:}
            \begin{itemize}
                \item Enhances critical thinking and self-awareness.
                \item Encourages continuous personal and professional growth.
            \end{itemize}
    \end{itemize}
\end{frame}

\begin{frame}[fragile]
    \frametitle{Key Points to Emphasize}
    \begin{itemize}
        \item Integrate Feedback Regularly: Feedback should be part of your routine.
        \item Be Open to Critique: Treat feedback as an opportunity for growth.
        \item Action on Feedback: Actively implement suggested changes.
    \end{itemize}
\end{frame}

\begin{frame}[fragile]
    \frametitle{Summary}
    Utilizing a variety of feedback mechanisms—Instructor feedback, peer reviews, and self-reflection—ensures a well-rounded approach to improving your capstone project. Each mechanism offers unique benefits, fostering an environment of collaboration, support, and continual learning.

    \textbf{Engagement with Feedback Methods:} Enhance the quality of your work and develop essential skills for your future academic and professional journey.
\end{frame}

\begin{frame}[fragile]
    \frametitle{Conclusion and Next Steps - Key Takeaways}
    \begin{itemize}
        \item \textbf{Project Planning:} Comprehensive planning with clear objectives, timelines, and deliverables is essential for success.
        
        \item \textbf{Research and Data Collection:} Use credible sources and methods for gathering high-quality data to enhance findings.
        
        \item \textbf{Execution and Iteration:} Implement your project in stages and refine your approach based on feedback.
        
        \item \textbf{Documentation:} Keep thorough records of each phase for presentation and future reference.
        
        \item \textbf{Presentation Skills:} Practice engaging storytelling to effectively communicate your findings.
    \end{itemize}
\end{frame}

\begin{frame}[fragile]
    \frametitle{Conclusion and Next Steps - Preparing Final Presentations}
    \begin{enumerate}
        \item \textbf{Review and Revise Presentation Materials:}
        \begin{itemize}
            \item Ensure clarity and visual engagement in slides.
            \item Use visuals such as graphs and charts to support your narrative.
        \end{itemize}
        
        \item \textbf{Practice Your Delivery:}
        \begin{itemize}
            \item Rehearse multiple times considering timing and clarity.
            \item Seek peer feedback for adjustments.
        \end{itemize}
        
        \item \textbf{Finalize Supporting Documents:}
        \begin{itemize}
            \item Prepare a comprehensive project report.
            \item Create handouts to enhance audience understanding.
        \end{itemize}
        
        \item \textbf{Prepare for Q\&A:}
        \begin{itemize}
            \item Anticipate questions and prepare responses.
            \item Practice answering questions with peers.
        \end{itemize}
    \end{enumerate}
\end{frame}

\begin{frame}[fragile]
    \frametitle{Conclusion and Next Steps - Importance of Feedback}
    \begin{itemize}
        \item \textbf{Incorporate Feedback:} Use insights from reviews and instructor feedback to improve your presentation quality.
        
        \item \textbf{Self-Reflection:} Critically assess your work to identify areas for enhancement before your presentation.
    \end{itemize}
    
    \textbf{Summary:} Completing your capstone project is a significant achievement. Focus on delivering a polished, engaging presentation that reflects your hard work and learning. Prepare thoroughly, be confident, and share your insights enthusiastically!
\end{frame}


\end{document}